\documentclass{article}
\usepackage{graphicx} % Required for inserting images
\usepackage{amsmath}

\title{CS 554 - Project Proposal}
\author{Alex Shaffer: aws9 }
\date{febuary 2026}

\begin{document}

\maketitle

\section*{Overview}
This project is an implementation and analysis project for a combination of two radial basis function (RBF) methods. 
These are the RBF partition of unity method (RBF-PU) and the RBF rational approximantion method (RBF-RA). 
These methods both provide oppertunities for parallel gains which will be explored and unitlized in this project.
The key idea of RBF methods is to represent a function or solution to a PDE as a linear combination of radial basis function centered on nodes.

\[
u(x) = \sum_i^N \lambda_i \varphi(||x-x_i||)
\]

\section*{Motivation}
RBF methods have some useful properties that make them desirable for interpolation and collocation methods.
Due to the Mairhuber-Curtis theorem (CITATION HERE), RBF methods are provably the only interpolation methods that can be used on arbitrary nodes in higher dimensions.
Additionally, the Gaussian RBF exhibits spectral convergence, meaning that the error decreases faster than any polynomial as the number of nodes increases.

\[
\text{Gaussian Kernel}: \varphi(r) = e^{-\epsilon r^2}
\]

There are considerable challenges to using RBF methods with the niave direct method.
The traditional approach suffers from ill-conditioning and high computational cost from dense linear solves as the number of nodes increases.

\subsection*{RBF-PU}
The RBF-PU method decomposes the domain into a set of overlapping patches and constructs local RBF interpolants on each patch. 
These patches are blended smoothly together using a partition of unity to create a global interpolant.
The global interpolant is given by:
\[
u(x) = \sum_{j=1}^M w_j(x) u_j(x)
\]
where $w_j(x)$ are the partition of unity weights and $u_j(x)$ are the local RBF interpolants on each patch.
The upshot is that for for patches that do not contain $x$, the corresponding weight $w_j(x)$ is zero. 
This adds a degree of sparsity to the system and allows for parallelization across patches.
A key step in solving PDEs with an iterative solver like GMRES is the application of the operator to a vector, which can be done in parallel across patches.
While the RBF-PU method helps the conditioning, due to improved conditioning on small patches, it does not fully resolve the issue of ill-conditioning as the number of nodes increases.

\subsection*{RBF-RA}
The RBF-RA method is a recent technique used for stabalization of RBF methods.
With the Gaussian RBF, spectral accuracy is achieved as the shape parameter $\epsilon$ approaches zero, but this also leads to ill-conditioning, because the RBFs become increasingly flat and linearly dependent.
The key idea of the RBF-RA method is to construct a rational approximation to the RBF interpolant, which can be evaluated stably even for small shape parameters.
This is done by solving many RBF interpolation problems for different values of $\epsilon$ on a safe contour in the complex plane.
A least squares problem is then solved to find the coefficients of the rational approximation, which can be evaluated stably for small $\epsilon$.

This method provides an oppertunity for parallelization across the different interpolation problems and the large least square problem, which has exploitable sparsity.



\end{document}